% ============================================================================
% chapters/ch1_exec.tex - Chapter 1: Executive Summary, CONOPS & RFP Matrix
% ============================================================================
\chapter{תקציר מנהלים, רקע המכרז וחוזה ההקמה הלאומי}
\label{ch:exec}

\section{פתח דבר: הדמוקרטיה הדליברטיבית בעידן הלוחמה התודעתית}
מערכות בחירות דמוקרטיות ומוסדות שלטון נציגים נשענים מאז ומעולם על תשתית אפיסטמית משותפת: יכולתם של אזרחים לקבל מידע אמין, לנהל דיון ציבורי חופשי בספירה הציבורית (\en{Public Sphere}), ולגבש הכרעה מושכלת בקלפי. בעשור האחרון חל שינוי מבני עמוק באופי הזירה הציבורית. השילוב בין אלגוריתמים של רשתות חברתיות המתגמלים זעם וקיטוב (\en{Engagement-Driven Amplification}), לבין זמינותם של מודלי שפה גדולים (\en{LLMs}) וכלי סינתזה מולטי-מודאליים (\en{Deepfakes}), העניק לגורמים עוינים --- מדינתיים ושאינם מדינתיים --- יכולת חסרת תקדים להנדס את השיח הלאומי בזמן אמת.

מבצעי השפעה זרים (\en{Foreign Information Manipulation and Interference -- FIMI}) אינם מסתפקים עוד בהפצת ידיעה כוזבת נקודתית. במקומם מופעלים ``נחילי תודעה'': אלפי חשבונות מזויפים בתיאום סינכרוני המניעים הסטה מלאכותית של סדר היום (\en{Agenda Hijacking}), מזריקים מסגרות פרשנות מקטבות (\en{Polarizing Frames}), ומפיצים הקלטות קוליות או סרטוני וידאו מזויפים בשעות הקריטיות טרם סגירת הקלפיות. במצב זה, מנגנוני התגובה המסורתיים של בדיקת עובדות בדיעבד (\en{Reactive Fact-Checking}) נכשלים מבחינה מתמטית ותקשורתית.

\begin{tenderbox}{ייעוד המסמך והצהרת כשירות למכרז הלאומי}
דוח מכין זה מוגש כמענה רשמי למכרז הבין-משרדי \textbf{\en{GOV-IL-RFP-2026-CSAI}} מטעם \textbf{ועדת הבחירות המרכזית}, \textbf{מערך הסייבר הלאומי} ו\textbf{משרד המשפטים}. מטרת המסמך להוכיח היתכנות אקדמית, הנדסית, כמותית וחוקתית להקמת מערכת \textbf{\en{CivicShield AI}} (``מגן התודעה הלאומי'') --- תשתית לאומית לחיסון תודעתי פרואקטיבי (\en{Cognitive Prebunking}) המבוססת על ועדת סוכני \en{AI} מתואמת, תחת פיקוח שיפוטי משולש השולל כל אפשרות לצנזורה אוטומטית.
\end{tenderbox}

\section{ניתוח הפער האסטרטגי: מבדיקת עובדות בדיעבד לחיסון פרואקטיבי}
הפרדיגמה הרווחת כיום בגופי רגולציה ובמערכות חדשות מבוססת על מעגל תגובה לינארי ואיטי: (א) ידיעה כוזבת עולה לרשת; (ב) היא צוברת חשיפה ויראלית במשך שעות; (ג) הציבור מתבצר בעמדות רגשיות; ו-(ד) רק לאחר שהנזק התרחש מתפרסמת הפרכה פוסט-פקטום. מערכת \en{CivicShield AI} מחליפה פרדיגמה זו בתיאוריית \textbf{החיסון התודעתי (\en{Psychological Inoculation / Prebunking})}, המאתרת דפוסי מניפולציה בעודם באיבם (\en{Weak Signals}) ומסייעת לגורמי ציבור מוסמכים לשדר ``נוגדנים קוגניטיביים'' ניטרליים טרם הגעת הנרטיב לזרם המרכזי.

\begin{table}[H]
\centering
\caption{השוואה אסטרטגית בין הגישה המסורתית לגישת \en{CivicShield AI}}
\label{tab:strategic_comparison}
\begin{tabularx}{\textwidth}{|>{\bfseries}p{3.2cm}|X|X|}
\hline
\rowcolor{csNavyBg}
ממד השוואה & \textbf{גישה מסורתית (בדיקת עובדות בדיעבד)} & \textbf{גישת \en{CivicShield AI} (חיסון פרואקטיבי)} \\
\hline
תזמון התערבות & לאחר ויראליות שיא (תגובה מאוחרת בשעות או ימים) & בזמן התהוות דפוס ההשפעה (התראה תוך שניות וחיסון מוקדם) \\
\hline
מוקד הניתוח & אימות עובדתי נקודתי של משפט בודד & זיהוי קמפיין מתואם (\en{CIB}), מסגור מפלג וטופולוגיית בוטים \\
\hline
מנגנון הגנה קוגניטיבי & תיקון שגיאה לאחר היווצרות הטיה רגשית & הקניית עמידות לטכניקת ההונאה (\en{Prebunking}) \\
\hline
מדרגיות הנדסית & תלויה בצוואר בקבוק של עורכים אנושיים & ועדת סוכני \en{AI} היררכית (\en{Multi-Agent Mesh}) + בקרה אנושית \\
\hline
מגבלה חוקתית & סיכון לשרירותיות רגולטורית בהסרת תוכן & איסור אכיפה ללא חתימה משולשת; הפרדה בין זיהוי לסמכות שיפוט \\
\hline
\end{tabularx}
\end{table}

\section{תפיסת ההפעלה המבצעית (\en{CONOPS})}
תפיסת ההפעלה הלאומית (\en{Concept of Operations -- CONOPS}) של מערכת \en{CivicShield AI} מתארת מחזור חיים סגור בין ארבעה שלבים רציפים הפועל 24/7 ומתגבר לקצב תגובה תחת \en{250ms} בתקופת 90 הימים שלפני בחירות ארציות:
\begin{enumerate}
  \item \textbf{איסוף ונירמול פומבי מאומת (\en{Privacy-Preserving Public Ingestion}):} צריכת מידע פומבי גלוי בלבד (\en{OSINT}) ללא גישה למסרים פרטיים, תוך מיסוך אנונימי תחת תקן פרטיות דיפרנציאלית ($\epsilon$-DP).
  \item \textbf{תזמור ניתוח רב-סוכני (\en{Multi-Agent Diagnostic Reasoning}):} הזרמת אירועים חריגים ל-\en{Agent Mesh Hub} שבו פועלים במקביל סוכן רשת ובוטים, סוכן מסגור תקשורתי, סוכן אימות פורנזי וסוכן תחזית אפידמיולוגית (\en{SIRI}).
  \item \textbf{שיפוט חוקתי ושער אדם (\en{Constitutional Gate \& Triple-Signature}):} בחינה דטרמיניסטית של חופש הביטוי, שיוך זר (\en{FIMI Attribution}) וסכנה לטוהר הבחירות; כל המלצת אכיפה דורשת חתימה קריפטוגרפית משולשת.
  \item \textbf{תגובה מאוזנת --- התראה והנגשת חיסון תודעתי:} הפקת דוחות שקיפות טכניים ומסרי \en{Prebunking} אובייקטיביים לתקשורת, לוועדת הבחירות ולציבור.
\end{enumerate}

\begin{figure}[H]
\centering
\begin{tikzhebrew}[
  node distance=0.85cm,
  stage/.style={rectangle, rounded corners=4pt, draw=#1, line width=1.2pt, fill=#1!8, text width=3.15cm, minimum height=1.85cm, align=center, font=\small},
  audit/.style={rectangle, rounded corners=4pt, draw=csNavy, dashed, line width=1pt, fill=csNavyBg, text width=11.2cm, minimum height=1.2cm, align=center, font=\small},
  arr/.style={-{Stealth[length=2.5mm]}, line width=1.1pt, color=csNavy}
]
  \node[stage=csNavy]    (s1) {\textbf{\heb{1. איסוף גלוי}}\\[3pt]\footnotesize\heb{טלגרם, רשתות, חדשות}\\[2pt]\texttt{\scriptsize OSINT + $\epsilon$-DP Filter}};
  \node[stage=csTeal,    right=of s1] (s2) {\textbf{\heb{2. ועדת סוכני \en{AI}}}\\[3pt]\footnotesize\heb{רשתות, מסגור, פורנזיקה}\\[2pt]\texttt{\scriptsize SIRI + ReAct Engine}};
  \node[stage=csCrimson, right=of s2] (s3) {\textbf{\heb{3. שער חוקתי}}\\[3pt]\footnotesize\heb{הפרדת סמכויות ובג"ץ}\\[2pt]\texttt{\scriptsize Triple-Key Human Sign}};
  \node[stage=csAmber,   right=of s3] (s4) {\textbf{\heb{4. תגובה ומגן}}\\[3pt]\footnotesize\heb{קמפיין חיסון והתראה}\\[2pt]\texttt{\scriptsize Prebunking + Audit}};

  \draw[arr] (s1) -- (s2);
  \draw[arr] (s2) -- (s3);
  \draw[arr] (s3) -- (s4);

  \node[audit, below=0.9cm of $(s2.south)!0.5!(s3.south)$] (aud) {
    \textbf{\heb{לולאת למידה וביקורת (\en{Audit Log}):}}\\[2pt]
    \footnotesize\heb{תיעוד קריפטוגרפי בלתי ניתן לשינוי וכיול מחודש של ספי ה-\en{SIRI} וזיהוי \en{Deepfake}}
  };

  \draw[arr, dashed] (s4.south) |- (aud.east);
  \draw[arr, dashed] (aud.north -| s2.south) -- (s2.south);
\end{tikzhebrew}
\caption{תפיסת ההפעלה המבצעית (\en{CONOPS}) מקצה לקצה של מערכת \en{CivicShield AI}}
\label{fig:conops}
\end{figure}

\section{מטריצת עמידה בדרישות המכרז הבין-משרדי (\en{RFP Compliance Matrix})}
\begin{table}[H]
\centering
\caption{מטריצת עמידה והתאמה לדרישות המכרז הממשלתי המשולב}
\label{tab:rfp_matrix}
\begin{tabularx}{\textwidth}{|>{\bfseries\en}p{1.8cm}|p{3.8cm}|X|p{2.0cm}|}
\hline
\rowcolor{csNavyBg}
\heb{\textbf{מזהה}} & \textbf{דרישת המכרז} & \textbf{מענה תכנוני ואדריכלי בדוח} & \textbf{עוגן בדוח} \\
\hline
RFP-01 & תיאום בין תקשורת פוליטית ומדעי המדינה להנדסת \en{AI} & תרגום פורמלי של 7 תיאוריות יסוד (קביעת סדר יום, מסגור, ספירלת השתיקה, חיסון) לרכיבי קוד וצינורות חישוב. & פרק 2 \\
\hline
RFP-02 & גילוי מוקדם של התנהגות מתואמת (\en{CIB}) ורשתות \en{FIMI} & ניתוח גרפים דינמי, אשכול סינכרוניזציה בזמן תגובה וייחוס מקוריות טקסטואלית מול בוטים. & פרקים 3, 4 \\
\hline
RFP-03 & זיהוי \en{Deepfake} מולטי-מודאלי בזמן אמת & מנוע איחוד בייסיאני (\en{Audio + Video + Metadata Provenance}) בעל $\text{AUC} > 0.96$. & פרק 5 \\
\hline
RFP-04 & חיזוי התפשטות נרטיבים וקיטוב אפקטיבי & הרצת משוואות דיפרנציאליות \en{SIRI} ומודל \en{Hegselmann-Krause} לכימות סף התפרצות ($R_0$). & פרק 5 \\
\hline
RFP-05 & זמינות גבוהה וביצועים ביום הבחירות & ארכיטקטורת \en{M/M/c} אלסטית, השהיה ממוצעת מתחת ל-\en{250ms} תחת שרידות \en{99.99\%}. & פרק 5 \\
\hline
RFP-06 & הגנה על חופש הביטוי ואיסור צנזורה אלגוריתמית & סוכן בקרה חוקתי מופרד ארגונית ומנגנון חתימה קריפטוגרפית משולש המחייב אישור שופט. & פרק 6 \\
\hline
RFP-07 & ישימות כלכלית, תוכנית עבודה וניהול סיכונים & תוכנית 4-שלבית ל-24 חודשים, ניתוח עלות-תועלת רב-שנתי ומערך מדדי \en{KPI/SLA}. & פרק 7 \\
\hline
RFP-08 & מתודולוגיית תזמור רב-סוכנים לפי עקרונות הקורס & הפעלת ועדת 5 סוכנים (\en{Master, Theory, Architect, Simulator, RedTeam}) ובקרת איכות רציפה. & פרק 8 \\
\hline
\end{tabularx}
\end{table}

\section{תרומת הפרויקט ומבנה הדוח המכין}
\begin{infobox}{מפת הדרכים של שמונת פרקי הדוח}
\begin{itemize}
  \item \textbf{פרק 2 (תשתית תיאורטית):} מנתח לעומק כיצד תקיפות מידע משבשות מנגנונים של \en{Agenda Setting}, \en{Framing} והמרחב הציבורי של הברמאס.
  \item \textbf{פרק 3 (אפיון דרישות ותרחישים):} ממפה 15 דרישות פונקציונליות (\en{FR}), 8 דרישות ביצוע (\en{NFR}) ו-4 תרחישי קיצון של יום הבחירות.
  \item \textbf{פרק 4 (ארכיטקטורה וועדת סוכנים):} פורש ארכיטקטורה בת 5 שכבות ואת מנגנון לולאת \en{ReAct}, לוח ה-\en{TODO} וה-\en{Hub} התקשורתי.
  \item \textbf{פרק 5 (מודלים מתמטיים וסימולציות):} מציג פיתוח מתמטי פורמלי וארבע הרצות סימולציה מדעיות ב-\en{Python}.
  \item \textbf{פרק 6 (משפט חוקתי ואתיקה):} מנסח את מסגרת הבלימה החוקתית, הפרטיות הדיפרנציאלית ומנגנון החתימה המשולש.
  \item \textbf{פרק 7 (סיכונים, מפת דרכים ותקציב):} מעמיד לוח גאנט דו-שנתי, תקציב מפורט וניתוח סיכוני אבטחה.
  \item \textbf{פרק 8 (מתודולוגיה וביבליוגרפיה):} מדגים כיצד הדוח עצמו נבנה ונערך באמצעות ועדת המומחים הרב-סוכנית לפי שיעור 4.
\end{itemize}
\end{infobox}
